\appendix

\section{Case studies}\label{app:casestudies}

\subsection{Lorenz-63}
\todo{Lorenz-63 setup (stochastic Lorenz-63 with wind forcing, parameter randomisation,
window length, regular and random observing systems, \Sz{}/\So{} model errors), to be
completed with the hand-ported L63 runs.}

\subsection{Lorenz-96}
\paragraph{Dynamics.}
The two-scale system has $K = 8$ slow variables $X_k$ and $J = 4$ fast variables
$Y_{j,k}$ per slow variable,
\begin{align*}
  \dot X_k &= X_{k-1}(X_{k+1} - X_{k-2}) - X_k + F - h \textstyle\sum_{j} w_j Y_{j,k} + c_1 W(t),\\
  \dot Y_{j,k} &= \big[Y_{j+1,k}(Y_{j-1,k} - Y_{j-2,k}) - Y_{j,k} + h_x X_k\big] / \epsilon ,
\end{align*}
with periodic indices, nominal values $F = 8$, $h = h_x = c_1 = 1$, $\epsilon = 0.1$ and
coupling weights $w = (1, 1, 0.1, 0.1)$, and a slow external forcing $W(t)$ (an AR(1)
process plus a sinusoid). Each window draws its own parameters: $F$, $c_1$, $h_x$,
$\epsilon$ and $w$ are drawn uniformly within $\pm 20\,\%$ of their nominal values
($h$ is fixed).
The system is integrated with a fourth-order Runge--Kutta scheme, $\Delta t = 0.001$,
over windows of 3000 steps.

\paragraph{Observations.}
The observed state has 24 variables: the 8 slow variables and the first two fast
variables of each slow variable; the 16 other fast variables are never observed. The
observation-error variance is $\mathbf{R} = 0.5\,I$. Two observing systems are used:
\begin{itemize}[leftmargin=*,nosep]
  \item \textbf{Regular:} all 24 variables observed at 30 equally spaced times per window
    (every 100 steps, first step included);
  \item \textbf{Random:} the same windows observed at 10 to 100 times at stratified
    random positions (first step always observed, at least one observation per DA
    sub-window), with 4 to 16 of the 16 observable fast variables observed per window.
\end{itemize}
Tables~\ref{tab:s0} and~\ref{tab:s1} and Figs.~\ref{fig:budget-s0}
and~\ref{fig:budget-s1} use the regular observing system; the random one is analysed in
Sect.~\ref{sec:q1-density}. All schemes assimilate exactly the same observations.

\paragraph{Data sets.}
The test set has 200 windows. The neural schemes are trained on 1000 windows and
validated on 100, with a random observing system redrawn at every batch (10 to 300
observation times, 4 to 16 observed fast variables) and fresh observation noise. The DA
hyper-parameters are selected on 50 separate validation windows.

\paragraph{Model error.}
\Sz{} and \So{} are described in Sect.~\ref{sec:testbeds}. At \So{}, the DA model has
its parameters biased by $+10\,\%$, receives a corrupted forcing and is truncated to the
24 observed variables; in the \So-full variant, it keeps the 40 variables.

\subsection{QG ocean}
\paragraph{Dynamics.}
The two-layer quasi-geostrophic ocean is solved on a $64 \times 64$ doubly periodic grid
over 1000\,km, with a 2\,h time step, $\beta = 1.5 \times 10^{-11}$\,m$^{-1}$\,s$^{-1}$
and a layer-depth ratio $\delta = 0.25$. Each window draws its deformation radius
(12.75--17.25\,km), upper-layer mean flow (0.0425--0.0575\,m\,s$^{-1}$) and bottom drag.
The ocean is forced by the wind-stress curl of a low-order chaotic atmosphere of coupled
gyrostats, projected on 12 Fourier modes ($|k| \le 2$), whose regime changes drive the
large-scale circulation; about one test window in five is calm.

\paragraph{Observations.}
Only the upper-layer streamfunction $\psi_1$ is observed, on 3 randomly drawn meridional
columns per day (4.7\,\% of the grid), each once per day, with a 5\,\% observation
error, mimicking along-track satellite altimetry. Scores are computed on 100 test
windows of 30 days, as the mean over the two streamfunction layers $\psi_1, \psi_2$.

\paragraph{Model error.}
At \So{}, the wind given to the DA model has a $+15\,\%$ amplitude bias, a random error of
25\,\% of each mode's standard deviation and a 50\,km position error; the deformation
radius is biased by $-10\,\%$ and the bottom drag by $-50\,\%$; the observation errors
grow to 15\,\% white plus 15\,\% correlated per pass; and the DA model runs on a
$32 \times 32$ grid, which does not resolve the deformation radius.

\paragraph{Data sets.}
The neural schemes are trained on 5000 windows drawn from the same distribution, with
observation-noise augmentation (5--15\,\% white, 0--15\,\% correlated per pass) and a
random number of observed columns per day. \todo{QG training details for D-Unet and SDA
(window length, density draw).}


\section{Operators and proofs}\label{app:proofs}

\paragraph{Gaussian operator.}
For $q = \mathcal{N}(m, v)$ (per element, standardised units) and $\xv_\tau = \tau
\xv_1 + (1-\tau)\xv_0$, the pair $(\xv_1, \xv_\tau)$ is jointly Gaussian with
$\Cov(\xv_1, \xv_\tau) = \tau v$ and $\mathrm{Var}(\xv_\tau) = \tau^2 v +
(1-\tau)^2$, so
\begin{equation}
  \Psi_q = m + \frac{\tau v}{\tau^2 v + (1-\tau)^2}\,(\xv_\tau - \tau m)
         = m + \frac{\snr\, v}{1 + \snr\, v}\,\Big(\frac{\xv_\tau}{\tau} - m\Big).
\end{equation}
Substituting $\xv_\tau^\star = \tau\xv_1^\star + (1-\tau)\xv_0$ and taking the expectation over $\xv_0$ gives property (iii).

\paragraph{The chain at every $\tau$.}
\todo{Write out: (a) $\Psi_p(\xv_\tau, C) = \E[\xv_1^\star \mid \xv_\tau, C]$ and
$\Psi_{p_S} = \E[\xv_1^\star \mid \xv_\tau, S] = \E[\Psi_p \mid \xv_\tau, S]$; (b)
the two orthogonality relations; (c) the condition on the scheme's randomness;
(d) strict propriety from injectivity of Gaussian smoothing; (e) the
$\tau \to 1$ expansion and the $d\snr$ integral (Verdú 2010).}

\paragraph{Mixture operator.}
\todo{The posterior-weighted average for an empirical or kernel-density output;
nearest-member snapping of the atomic form at large $\tau$.}

\section{The design matrix}\label{app:matrix}
\todo{Five choices, each a column of Table~\ref{tab:schemes}:
F1 prior--likelihood split (factorised with explicit $\Hop$, $\mathbf{R}$ vs amortised);
F2 dynamics structure (one-step transition, joint window prior, none; equations vs samples);
F3 sequential inference (sequential vs whole window; exact only if F2 holds in $\xv$ and observation errors are white);
U-state (none; explicit $\mathbf{B}$, $\mathbf{P}$ or ensemble; implicit, learned);
U-model (none; state-space $\mathbf{Q}$ or inflation; parameter spread; structure).
Every scheme's row.}

\section{Tuning}\label{app:tuning}
\todo{Inflation sweeps of the ETKF, EnKF and ETKS (30 and 100 members) on the
validation windows, with the selection rule; Weak-4DVar model-error scale
$\times$ L-BFGS budget (Table from \texttt{docs/results/l96\_weak4dvar.md}); SDA
guidance weight; D-Unet+SDA start time and guidance.}

\section{Sensitivities}
\todo{Ensemble size; L-BFGS iterations; training budget and data; flow sampler
steps; the 400-epoch rows.}

\section{The fixed-observing-system protocol and network size}\label{app:tiers}
\todo{From \texttt{p1\_l96\_benchmark.md} annex A: S+ / M / L tiers, VanillaCFM vs
PredictStateCFM, $\tau = 0$ heads, SDA conditioning, under the earlier
fixed-observing-system protocol; not comparable with the main tables.}

\section{Estimating \texorpdfstring{$\FMS_\tau$}{FMS}}\label{app:fms}
\todo{Sliced vs per-element evaluation on L96 (snapshot and space--time random projections, 64 directions): same rankings and S0/S1 conclusions; the advantage of SDA over the ETKS at \Sz{} shrinks to 5--8\,\% on space--time patches.}
\todo{Monte Carlo over $\xv_0$ with common random numbers; kernel bandwidth;
24-channel snapshot blocks and their effective-sample-size degeneration at
$\tau = 0.75$; window bootstrap.}

\section{Rank histograms}
\todo{From \texttt{docs/results/l96\_rank\_histograms\_c4.md}.}
